\chapter[Exogenous Chemical Effects]{Exogenous Chemical Effects and the Stimulation of Germination by Gibberelins}

This chapter is divided into four parts: Stimulation of Germination by
Gibberelins, Parastic Plants, Stimulation of Germination by Nitrogen Compounds, and
Stimulation by Anaerobic Conditions.

\section[Gibberelins]{Stimulation of Germination by Gibberelins}

The work reported herein
and two recent papers (refs. 9 and 10) leave no doubt that gibberelins and particularly
Gibberelic Acid-3 (GA-3) will from now on play an important role in growing plants from
seeds. The reaction of the seeds of various species to GA-3 varies from dramatic
stimulation of germination to little effect to killing the seeds. Among the species
exhibiting dramatic stimulation of germination, the effect varies from producing normal
healthy seedlings to producing seedlings so badly etiolated that they soon die.
Concentrations and exposures can be varied, and it is possible with some species to
optimize.these variables and get healthy seedlings. This was found for Gentiana
verna (ref. 9, confirmed in present work) and for Linum capitatum and Ribes cereum. It
has been reported that cytokinins are helpful in reducing the etiolation (ref. 9).
However, there is no certainty that satisfactory conditions can be found with all of the
species that show stimulation, and most growers \textit{will prefer to avoid the use of GA-3}
and other gibberelins unless necessary such as in the following circumstance.

\textbf{It is now proposed that gibberelins are natural stimulators of
germination In certain species and that the evolution of this requirement
is critical to the survival of those species.} The recognition of this came about
when it was found that germination in Echinocereus pectinatus hybrids was less than
1\% and germination in the rosulate Violas was zero without the use of GA-3. Both of
these species are relatively tiny plants growing in harsh dry environments, The
Echinocereus from the dry windswept deserts of Southwestern United States and the
rosulate Violas from high in the dry windswept Chilean Andes. For these tiny
seedlings with their tiny root systems to survive, the seeds must fall into crevices in
rocks or deep into coarse gravel where there is a small pocket of moist leaf mold. The
fungal action in this leaf mold produces gibberelins which trigger the germination. In
this way germination is triggered in the only place where the seedlings could survive.
Some other species that give significant germination only with GA-3 are Evodia
danielli, some Gentiana, Physoplexis comosum, Ranunculus lyallii, Ribes cereum,
Romneya coulteri, Rubus parviflorum, most Thalictrum, Trillidium, and most Trillium.

A brief personal note may interest the reader. Claude Barr (author of \textit{Jewels of
the Plains}) had always cautioned me that Cactus seedlings cannot take direct sun and
dessicating conditions and that he had killed many cactus seedlings before he
recognized this. It now becomes clear how they have developed a clever detection
system for germinating in the one place in which they could survive.
The dosages and times of exposure can be regulated. For example with the
three species studied in detail (Gentiana verna, Linum capitatum, and Ribes cereum),
one-tenth as much (0.1 cubic millimeter) of GA-3 can be .used, and the seeds can be
removed each day, washed, and placed in clean moist towels. While all of this may be
somewhat qualitative and unscientific, the usual procedure of making up exact
concentrations and soaking the seeds for exact time periods is not as precise as might
be thought. There is selective absorption of the Gibberelin on the seed coats which
alters the effective concentrations, and the solutions cannot really be used again.

\textit{A totally unexpected result} was that most of the species that required light to
germinate would germinate in the dark if treated with GA-3. This obviously has
important implications in regard to the chemistry involved in the stimulation of
germination by both light and GA-3. It would be tempting to speculate, but for the
moment the examples are all documented in Chapter 20 under each species, and it is
best left at that.

About half of all Lilium are two step germinators in that they form a small bulb
and root at 70, but then need about three months at 40 before the true leaf will form on
returning to 70. Three of these (L. canadense, L. martagon, and L. szovitsianum) were
subjected to treatment with GA-3. The first step was unaffected both in respect to the
rate characteristics of the germination and the appearance of the seedlings. However,.
about three weeks after the initial germination and after the bulb and root had formed,
the true leaf emerged in about half of the L canadense. This effect is of some value
because it much reduces the.time required to get a photosynthesizing leaf.

Trillium are also two step germinators. The GA-3 treatment was tried on several
species as described in Chapter 20. With Trillium the rates of both steps increased
with the GA-3 treatment, and most Trillium germinated only when treated with GA-3.
The results with Shortia galacifolia are now described in detail because they
illustrate the intricate complexities and interactions. In the first edition it was reported
that Shortia galacifolia required New Jersey Pine Barren Sand plus light in order to
germinate (discovered by John Gyer). This was clearly an example of an exogenous
chemical. requirement for germination and was only the second such example at the
time (the other was Striga, discussed further on). Further work has shown that Shortia
is a particularly interesting example. It is now found that GA-3 stimulates the
germination in \textit{either} 70L or 70D whereas the New Jersey Pine Barren Sand
stimulated the germination \textit{only} in 70L. Clearly the chemical in the Sand is not GA-3,
although it is likely that it is a fungal product of the gibberelin type.

When sown on the surface of wet New Jersey Pine Barren Sand under the
fluorescent lights, germination was 60-70\% at a rate of 6-7\% per day with an induction
time of eight days. When a single layer of paper towel was placed between the seed
and the sand, germination occurred at the same rate and percentage. When four
layers of towel were placed between the seed and the sand, germination was reduced
to 30\%. The induction time was the same, and the rate was the same providing it was
calculated on the basis of the reduced number of seeds germinating instead of the
total number of seeds. When the seed was sown on the sand in the dark and shifted to
light after seventeen days, germination was only 20\% showing that a time in the dark
as short as seventeen days caused severe injury. The above treatments were the only
ones that gave germination.

The following procedures failed to give a \textit{single} germination. The seed was
placed on moist towels in dark and light. The-seeds were placed on towels in light,
and the towels were moistened with 1\% aqueous acetic acid, 5\% aqueous acetic acid,
or 1\% aqueous iron sequestrene. The seed was sown in light on sphagnum peat.
The seed was shifted from the towels to the sand after two weeks. The seed was dry
stored at 70 for four weeks and then sown on the New Jersey sand in light.
There is another situation that could be interpreted as an exogenous chemical
requirement. In growing orchids from seed in agar gel, many components are needed
in the gel, and this is discussed under Orchidaceae in Chapter 21. However, it is not
clear whether these exogenous chemical requirements are needed for the
germination itself or for the development of the seedling. The two aspects become
intertwined in Orchidaceae where the seed is so small and carries barely any store of
nutrient.

One is hesitant to add an anecdotal observation, but it relates to Fraxinus.and
Liriodendron, which have been singularly difficult to germinate, and it relates to fungal
action and the possibility of a gibberelin requirement for germination. A number of
truck loads of leaves and chopped branches were placed on our property to make a
great compost pile. Three years later there appeared a dozen seedlings of Fraxinus
americana and about one hundred seedlings of Liriodendron tulipifera. No
germination of these two species had occurred in the towel experiments. Related to
this is the report that embedding Crataegus seed in peat speeds up germination (ref.
22).

Before leaving this subject, the two most useful papers (refs. 9 and 10) will be
summarized, and the history of gibberelins briefly reviewed. Stewart and Presley (ref.
9) found marked stimulation of germination in Aquilegia jonesii x saximontana, several
Campanula, several Gentiana, Lewisia tweedyi, and Primula parryi. Nikolaeva and
coworkers (ref. 10) reported stimulation in Campanula latifolia, Lunaria rediviva,
Polygonum bistorta, Primula veris, and ten species of Trollius. Ten of these Trollius
germ. in 5-20 days if the seeds were soaked for two days in 500-1000 mg/I (equivalent
to ppm) of GA-3, and Trollius asiaticus germ. in 30-50 d after treatment. Normally
these species require 3-5 months of chilling before they will germinate in warmth, and
Trollius asiaticus require eleven months of chilling. In addition Prof. Roger Koide at
Pennsylvania State University found that germination of dry stored seed of Plantago
erecta was increased from the 2-5\% range to 80\% by one hour immersion of the seed
in 300 ppm aqueous solution of gibberelic acid.

The history of gibberelins is exhaustively reviewed in \textit{Gibberelins} (ref. 23).
Briefly, a Japanese chemist in 1926 showed that the "crazy rice" disease was caused
by a chemical secreted by a Gibberelus fungus growing on the rice. This chemical
was named gibberelin after the fungus from which it was derived. In 1936 another
Japanese chemist crystallized the active principle. In the years that followed many
chemicals were tested for growth promoting activity by applying a solution of the
compound to be tested in a block of gel to the side of an oat coleoptile. The degree of
bending and the concentration of the compound defined a quantitative measure of
activity. By the time the book \textit{Gibberelins} was published, 67 different gibberelins had
been isolated and their chemical structure proven.

My information at present is that Gibberelic Acid-3 (GA-3) is marketed by
several companies. They buy it from Abbott Laboratories, North Chicago, Illinois, who
are the actual producer. Generally orders for GA-3 will not be honored unless
originating from an academic institution or some formally recognized research group
so that it is not available to the general public.

Up to two years ago, although there was a general consensus in the Iliterature
that gibberelins were the most promising group for stimulating germination of seeds,
the reports were so fragmentary that the subject was only briefly mentioned in the first
edition of this book and the book \textit{Gibberelins} (ref. 23) did not have any Chapter or
hardly any mention of this subject. All of this changed with the two recent publications
(refs. 9 and 10).

Related to the initiation of germination by gibberelins is the blockage of
germination by abscicic acid. This is elegantly reviewed in the book \textit{Abscicic Acid}
(ref. 24). It is tempting to regard the whole subject of promotion and blockage of
germination as due to the ratio of gibberelin to abscicic acid, but I believe the matter to
be far more complex. Incidentally auxins did not stimulate germination (ref. 25).

\section{Parasitic Plants}

Although the literature gives the general impression that
parasitic and semiparasitic plants require proximity to the host for germination, this has
not been found to be true in the present work. Aureolaria, several Castillejas, and
Phoradendron (Mistletoe) all germinated on paper towels moistened only with water.
The Aureolaria and the Castillejas developed the two cotyledons and several true
leaves before dieing off because of lack of host. The Phoradendron developed only a
2-4 mm. green leaf like shoot. A dramatic exception to the above comes from the work
of David Lynn and others (reviewed in ref. 26) where it was unequivocally
demonstrated that Striga asiatica requires a specific hydroquinone (exuded from the
host) in order for germination to occur.

Striga asiatica is a serious parasite of grasses. It has been much studied
because it reduces grain production in Asia and there is the potential for this problem
to become serious in this country. Most of the studies have been conducted with
sorghum. Briefly, sorghum exudes droplets along the entire length of the roots. These
droplets are 95\% a single chemical whose-structure has been determined. Suffice to
say that it is a hydroquinone (HO) with a structure resembling phytol. This HQ is
readily oxidized to the quinone form Q which is inactive. The droplets of exudate
contain an oxidation inhibitor as a minor constituent and this inhibitor does slows
down the oxidative destruction of the HO. The inhibitor resembles the well known
antioxidant di-t-butyl-p-cresol and undoubtedly acts in the manner well documented
for the cresol. Even with the oxidation inhibitor, the HO does not survive beyond 5-10
mm. from the sorghum root so that the seeds of Striga must be that close to the
sorghum root for germination to occur. Incidentally, the HO is also required for
formation of the haustoria, and the seedling has just five days to develop the haustoria
and attach it to the host or the seedlings dies. It is possible that there are further levels
of response by the seedlings to the chemicals exuded from the host so that there is
much complexity.

Other chemical compounds can initiate the germination of Striga. Among these
are ethylene, gibberelins,cytokinins, NaOCI, and sulfuric acid. It would appear that the
initiation of germination is the result of the destruction of some specific germination
inhibitor. Most interesting is the fact that cotton secrets a substance similar to HQ
named strigol which also initiates germination of Striga, but cotton is not a suitable
host. Thus cotton can be planted in a field infested with Striga seed as a trap plant. It
initiates germination, but the seedlings perish and the infestation is reduced. The
seeds of Striga can survive over twenty years in soils and remain viable so that
destruction of Striga seed is important agriculturally.

\section[Nitrogen Compounds]{Stimulation of Germination by Nitrogen Compounds}

There have been
many fragmentary reports that compounds of nitrogen stimulate germination. The
most quantative study was on germination of red rice (refs. 2-7, 27, and 28) where it
was shown that nitrite stimulated germination and the stimulation was measured as
function of the pH. The curve resembles the the curve for the equilbrium between
nitrous acid and nitrite anion suggesting that nitrous acid is attacking some inhibitor.

This speculation would also explain various reports that nitrate; hydroxylamine,
ammonia, and thiourea (refs. 2-7, 27, and 28) sometimes stimulate germination. All
three are well known to be converted to nitrite in soils, nitrate by reduction and the
other two by oxidation.

Related to the question of exogenous chemical effects are situations like the
black walnut whose roots secrete a substance, juglone, which inhibits the growth of
many species. This could give the appearance of inhibiting germination, although the
action might be described as killing the seedling as the radicle emerges. It would be
difficult to distinguish the two effects, and in fact it merges into a semantic argument.
Seedlings have been raised from excised embryos, and some Iris and Lilium
hybrids require such treatment. This field has been omitted in this book, but obviously
the excised embryos need exogenous nutrient and perhaps other materials.

\section[Anaerobic Conditions]{Stimulation by Anaerobic Conditions}

Several reports in the literature
can be interpreted to suggest that germinations of Nymphaea, Typha latifolia, and
other aquatic species are stimulated by mildly anaerobic conditions. Germination of
Nymphaea and Typha has been reexamined. It was found that Nymphaea tetragona
germinated best in a 40-70 pattern and Typha latifolia required light for germination,
both in aerobic conditions. However, a mild stimulation of germination was observed
in Nymphaea tetragona and in Butomus umbellatus using mildly anaerobic conditions.

Seeds of Butomus umbellatus, Nymphaea tetragona, and Typha latifolia were
placed in the moist towels in a polyethylene bag. The bag was blown full by my
exhaling into it several times. The bag was collapsed and the processes repeated two
more times. In this simple way mildly anaerobic atmosphere was created in which
the oxygen was reduced and the carbon dioxide increased. Germination was mildly
stimulated at 70 for both Butomus umbellatus and Nymphaea tetragona but not Typha.
My conclusion is that mildly anaerobic conditions may play some role in the
germination of aquatic species under natural conditions, but it is not an absolute
requirement. For the pragmatic horticulturist aerobic conditions can be found that give
better results, at least for the three species cited above.

The literature reports are now briefly reviewed; Typha latifolia was reported to
germinate only under water, and it was proposed that the reduction in oxygen
pressure induced the germination (ref. 29). Nymphaea germination has been studied
by two groups. Else and Riemer found that crowding of Nymphaea odorata seeds
stimulated germination and that aeration inhibited germination. They attributed these
effects to ethylene formation (ref. 30). Roberta and Fred Case of Saginaw, Michigan,
also found that crowding of seeds of Nymphae stimulated germination. An effective
procedure was to place 1-2 inches of mud in a small glass jar, place the seeds on the
mud, add 2-3 inches of water, and close with a screw cap. Of course some oxygen is
required for the germination metabolism, so that oxygen cannot be completely
excluded, but only reduced.

A number of aquatic plants have been reported to require daily oscillations of
water temperature to achieve germination (ref. 31). It is possible that the temperature
oscillations develop mildly anaerobic conditions. As the temperature rises oxygen is
forced out of the water because the solubility of oxygen is lower at the higher
temperature; When the temperature drops the oxygen is slow to reestablish the
equilibrium solubility at the lower temperature so that mild anaerobic conditions
develop. Finally, the cocklebur (Xanthium canadense) contains two seeds in each
capsule. The upper one germinates in completely aerated media. The lower seed
require partially anaerobic conditions for germination (refs. 31 and 32).

